\chapter{Introducción}\label{introducciuxf3n}

\begin{center}\rule{0.5\linewidth}{0.5pt}\end{center}

\section{Planteamiento del Problema}\label{planteamiento-del-problema}

Las pequeñas y medianas empresas (pymes) constituyen el motor económico de América Latina. Según la Comisión Económica para América Latina y el Caribe (CEPAL, 2022), las pymes representan más del 99\% del tejido empresarial en la región y generan entre el 60\% y el 70\% del empleo formal. Sin embargo, pese a su relevancia, enfrentan una brecha estructural frente a las grandes corporaciones en materia de gestión operativa y adopción tecnológica.

Un diagnóstico recurrente en la literatura especializada señala que la gran mayoría de las pymes latinoamericanas aún gestionan sus operaciones ---inventario, ventas, clientes y contabilidad básica--- mediante herramientas rudimentarias: cuadernos físicos, hojas de cálculo de Excel no estructuradas o sistemas ad hoc sin integración entre módulos. De acuerdo con el Banco Interamericano de Desarrollo (BID, 2021), solo el 16\% de las pymes de la región ha adoptado algún tipo de software de gestión empresarial, frente a un 74\% en países de la OCDE.

Esta brecha de digitalización se traduce en consecuencias concretas y medibles para el negocio:

\begin{itemize}
\tightlist
\item
  \textbf{Errores de inventario}: la gestión manual provoca discrepancias frecuentes entre el stock físico y el registrado, generando roturas de stock o sobrestock. Ambos escenarios tienen costos directos: ventas perdidas en el primer caso, capital inmovilizado en el segundo.
\item
  \textbf{Ineficiencia operativa}: los empleados dedican tiempo significativo a tareas administrativas repetitivas ---conciliación de registros, elaboración de reportes manuales, búsqueda de información--- tiempo que podría orientarse a actividades de mayor valor.
\item
  \textbf{Toma de decisiones reactiva}: sin acceso a datos consolidados y en tiempo real, los dueños de pymes toman decisiones basadas en intuición antes que en evidencia. El abastecimiento de productos, la identificación de clientes de alto valor y la detección de tendencias de ventas se realizan de manera subjetiva, sin soporte de indicadores ni reportes estructurados.
\end{itemize}

Las soluciones existentes en el mercado presentan limitaciones específicas para este segmento. Los sistemas ERP de nivel empresarial (SAP, Oracle) son prohibitivos en costo y complejidad para una pyme. Las alternativas de código abierto como Odoo, si bien potentes, requieren infraestructura propia y capacidad técnica para su instalación, configuración y mantenimiento. Las soluciones SaaS disponibles en el mercado latinoamericano (Alegra, Siigo, Defontana) se enfocan principalmente en facturación electrónica y contabilidad, sin integrar capacidades de gestión de inventario, análisis de clientes y reportes operativos en un único entorno accesible.

Se identifica así un problema de investigación y desarrollo bien definido: \textbf{la inexistencia de una plataforma SaaS integrada, accesible y orientada al contexto latinoamericano, que combine la gestión operativa esencial de una pyme ---inventario, ventas, clientes y reportes--- bajo un modelo de precios y complejidad adecuados para este segmento.}

\section{Justificación}\label{justificaciuxf3n}

\subsection{Justificación Económica y Social}\label{justificaciuxf3n-econuxf3mica-y-social}

La digitalización de las pymes no es un objetivo exclusivamente tecnológico; es un imperativo económico. El BID estima que cerrar la brecha digital de las pymes en América Latina podría incrementar el PIB regional en hasta un 5\% adicional para 2030. En términos individuales, los estudios de adopción de ERP en pymes reportan reducciones de costos operativos del 15\% al 30\% y mejoras en la eficiencia de la gestión de inventario de hasta el 40\% en el primer año de implementación (Kumar \& van Hillegersberg, 2000; Duan et al., 2012).

La relevancia del problema es, por tanto, tanto académica como práctica: una solución bien diseñada tiene potencial de impacto directo en la competitividad de cientos de pequeñas empresas.

\subsection{Justificación Tecnológica}\label{justificaciuxf3n-tecnoluxf3gica}

Las decisiones tecnológicas que sustentan OrbitEngine no se justifican por la mera disponibilidad de las herramientas seleccionadas, sino por su capacidad de reducir el riesgo técnico del proyecto y acelerar su tiempo de salida al mercado en un contexto de equipo reducido y plazos académicos acotados:

\begin{enumerate}
\def\labelenumi{\arabic{enumi}.}
\tightlist
\item
  \textbf{Idoneidad del stack para una plataforma transaccional multi-tenant}: la elección de un backend asíncrono basado en FastAPI con tipado estricto, junto con un frontend declarativo en React, responde a requisitos concretos del dominio ---validación rigurosa de datos contables, concurrencia en operaciones de inventario y trazabilidad por organización---, y no a una preferencia genérica por tecnologías populares.
\item
  \textbf{Contrato API como única fuente de verdad}: la adopción de un enfoque \emph{contract-first} mediante OpenAPI permite generar automáticamente el cliente tipado del frontend a partir del esquema del backend, garantizando consistencia entre capas, detectando regresiones en tiempo de compilación y eliminando una clase entera de errores de integración propios de los desarrollos manuales.
\item
  \textbf{Seguridad desde el diseño}: la arquitectura incorpora autenticación basada en JWT, control de acceso por roles a nivel de endpoint y aislamiento lógico de datos por organización como invariantes del modelo, no como capas añadidas a posteriori, lo cual es crítico al manejar información financiera de terceros.
\item
  \textbf{Aprovechamiento de inteligencia artificial generativa en el ciclo de desarrollo}: la incorporación de asistentes de IA en tareas de codificación, documentación y revisión permitió al equipo sostener un ritmo de entrega comparable al de equipos de mayor tamaño, manteniendo estándares de calidad mediante revisión humana sistemática del código generado.
\item
  \textbf{Viabilidad económica del despliegue en la nube para el segmento pyme}: las plataformas de despliegue gestionado seleccionadas ofrecen, en 2026, un costo marginal por organización lo suficientemente bajo como para que el modelo de suscripción resulte rentable incluso con un volumen reducido de clientes, condición que no se cumplía con la infraestructura disponible una década atrás.
\end{enumerate}

\subsection{Justificación Académica}\label{justificaciuxf3n-acaduxe9mica}

La pertinencia académica del presente proyecto se sustenta en tres ejes complementarios que articulan su contribución al conocimiento, su rigor metodológico y su correspondencia con la formación recibida en el programa.

En primer lugar, el proyecto aporta evidencia empírica a un vacío específico de la literatura: si bien la adopción de sistemas de información en pequeñas y medianas empresas ha sido ampliamente documentada en contextos anglosajones y europeos (Kumar \& van Hillegersberg, 2000; Duan et al., 2012), los estudios centrados en pymes latinoamericanas ---y en particular sobre plataformas SaaS multi-tenant diseñadas para sus restricciones de costo, capacidad técnica y conectividad--- son notablemente más escasos. OrbitEngine se propone como un caso de estudio que documenta tanto el diseño arquitectónico como los resultados de su adopción en este contexto regional.

En segundo lugar, la validación del sistema sigue un diseño cuasi-experimental de tipo pre/post con empresas piloto, articulando métricas cuantitativas ---reducción del tiempo en tareas administrativas, disminución de la tasa de error en operaciones de inventario y tiempos de respuesta en producción--- con métricas cualitativas estandarizadas, en particular el instrumento \emph{System Usability Scale} (SUS) y entrevistas semiestructuradas de cierre. Este enfoque mixto permite contrastar con evidencia las hipótesis sobre el impacto de la digitalización en la eficiencia operativa, más allá de la sola demostración funcional del software.

En tercer lugar, el proyecto se inscribe en la línea de trabajo del semillero \emph{Software como Innovación} y articula de manera integral las competencias formativas del pregrado en Ciencias de la Computación e Inteligencia Artificial: ingeniería de requisitos, arquitectura de software, desarrollo full-stack, despliegue en la nube, validación experimental y comunicación de resultados. Lo anterior justifica su carácter como proyecto de grado y no como un desarrollo de software meramente aplicado.

\section{Objetivos}\label{objetivos}

\subsection{Objetivo General}\label{objetivo-general}

Desarrollar e implementar una plataforma SaaS multi-tenant para la gestión integral de procesos internos en pequeñas y medianas empresas, que integre módulos de gestión de inventario, ventas, clientes y reportes operativos, y validar su impacto en la eficiencia operativa mediante pruebas con empresas reales.

\subsection{Objetivos Específicos}\label{objetivos-especuxedficos}

\begin{enumerate}
\def\labelenumi{\arabic{enumi}.}
\item
  \textbf{Diseñar} la arquitectura técnica de una plataforma SaaS multi-tenant que garantice el aislamiento de datos entre organizaciones, la escalabilidad horizontal y la seguridad de la información, documentando las decisiones arquitectónicas con su justificación.
\item
  \textbf{Desarrollar} los módulos de gestión operativa esenciales ---autenticación y control de acceso basado en roles, gestión de inventario con alertas automáticas, registro de ventas con generación de documentos, y gestión de clientes con análisis de historial--- conforme a los requisitos levantados con usuarios reales de pymes.
\item
  \textbf{Construir} un sistema de reportes y analítica que proporcione indicadores clave de desempeño (KPIs) en tiempo real en el dashboard, junto con la capacidad de exportar a Excel los listados operativos (inventario, ventas y clientes), facilitando la toma de decisiones basada en datos.
\item
  \textbf{Desplegar} la plataforma en infraestructura de nube (Railway para el backend y la base de datos, Vercel para el frontend) con un pipeline de integración y entrega continua (CI/CD), garantizando una disponibilidad mínima del 95\% y tiempos de respuesta inferiores a 2 segundos para operaciones transaccionales.
\item
  \textbf{Validar} la solución mediante pruebas de usabilidad y rendimiento con al menos dos empresas piloto, midiendo el impacto en la eficiencia operativa a través de métricas cuantitativas (reducción de tiempo en tareas, tasa de error) y cualitativas (satisfacción de usuario).
\end{enumerate}

\section{Alcance y Limitaciones}\label{alcance-y-limitaciones}

\subsection{Alcance del Sistema}\label{alcance-del-sistema}

OrbitEngine abarca el siguiente conjunto de funcionalidades dentro del producto mínimo viable (MVP):

\textbf{Módulos implementados:}

\begin{itemize}
\tightlist
\item
  Autenticación y gestión de usuarios con control de acceso basado en roles (RBAC): Administrador, Vendedor y Visualizador.
\item
  Gestión de inventario: CRUD de productos por categorías, control de stock en tiempo real, alertas automáticas de nivel mínimo e historial de movimientos.
\item
  Gestión de ventas: registro de transacciones multi-producto con generación de número de factura, historial de ventas con filtros y búsqueda.
\item
  Gestión de clientes: base de datos de clientes, historial de compras y métricas de comportamiento.
\item
  Dashboard y reportes: KPIs en tiempo real (ventas del día/mes, stock bajo, top de productos, tendencia de ventas) y exportación a Excel de los listados de inventario, ventas y clientes.
\item
  Infraestructura multi-tenant: soporte para múltiples organizaciones con aislamiento total de datos.
\end{itemize}

\textbf{Características técnicas:}

\begin{itemize}
\tightlist
\item
  API RESTful documentada con OpenAPI/Swagger.
\item
  Interfaz web responsive, compatible con navegadores modernos (Chrome, Firefox, Safari, Edge).
\end{itemize}

\textbf{\emph{- Despliegue en Railway (backend) y Vercel (frontend) con GitHub Actions.}}

\begin{itemize}
\tightlist
\item
  Suite de pruebas automatizadas con cobertura mínima del 60\%.
\end{itemize}

\subsection{Limitaciones y Exclusiones}\label{limitaciones-y-exclusiones}

Las siguientes características quedan fuera del alcance del MVP y podrán abordarse en versiones futuras:

\begin{itemize}
\tightlist
\item
  Facturación electrónica oficial con validez tributaria (DIAN, SAT, SII u organismos equivalentes).
\item
  Gestión de nómina y recursos humanos.
\item
  Aplicación móvil nativa (iOS/Android).
\item
  Integración con pasarelas de pago o sistemas ERP externos.
\item
  Soporte para gestión de múltiples sucursales o bodegas dentro de una misma organización.
\item
  Gestión de servicios (el sistema está orientado a empresas de comercio de productos físicos).
\item
  Múltiples monedas o idiomas en la interfaz.
\end{itemize}

\subsection{Supuestos}\label{supuestos}

\begin{enumerate}
\def\labelenumi{\arabic{enumi}.}
\tightlist
\item
  Los usuarios de la plataforma disponen de conexión a internet estable para el acceso a la aplicación web.
\item
  Las empresas piloto están dispuestas a cargar sus datos de productos, clientes y stock inicial para comenzar a operar con el sistema.
\item
  El equipo de desarrollo cuenta con cuentas activas en Railway y Vercel con planes adecuados para el despliegue de la plataforma.
\end{enumerate}

\section{Metodología}\label{metodologuxeda}

El proyecto se desarrolló bajo un enfoque de investigación aplicada con componente de desarrollo de software. Se adoptó la metodología ágil Scrum adaptada para equipos académicos, organizando el desarrollo en sprints de dos semanas durante un período de siete meses (octubre 2025 -- abril 2026).

La investigación siguió las siguientes etapas:

\begin{enumerate}
\def\labelenumi{\arabic{enumi}.}
\tightlist
\item
  \textbf{Documentación y planificación} (octubre 2025): revisión bibliográfica, levantamiento de requisitos con usuarios potenciales, diseño de arquitectura, modelo de datos y definición del backlog inicial.
\item
  \textbf{Desarrollo del núcleo} (noviembre 2025 -- enero 2026): implementación de la infraestructura base y los módulos de autenticación, inventario, ventas, clientes y reportes.
\item
  \textbf{Estabilización y refinamiento} (febrero -- marzo 2026): corrección de errores, mejoras de experiencia de usuario, pruebas de carga y consolidación del despliegue en producción.
\item
  \textbf{Validación con empresas piloto} (abril 2026, semanas 1--3): pruebas con usuarios reales, recolección de métricas de eficiencia operativa y encuestas de usabilidad.
\item
  \textbf{Documentación y entrega} (abril 2026, semana 4): consolidación del informe final, análisis de resultados y defensa del proyecto.
\end{enumerate}

\section{Estructura del Documento}\label{estructura-del-documento}

El presente informe se organiza en los siguientes capítulos:

\begin{itemize}
\tightlist
\item
  \textbf{Capítulo 2 --- Marco de Referencia}: define los conceptos fundamentales del dominio (pyme, ERP, SaaS, multi-tenancy, API REST, RBAC) y presenta el estado del arte mediante una revisión de soluciones existentes en el mercado y de la literatura académica relevante.
\item
  \textbf{Capítulo 3 --- Análisis y Diseño del Sistema}: describe el proceso de levantamiento de requisitos, la arquitectura técnica adoptada, el diseño del modelo de datos y el diseño de las interfaces de usuario.
\item
  \textbf{Capítulo 4 --- Desarrollo e Implementación}: detalla la metodología de desarrollo utilizada, los módulos implementados, las decisiones técnicas relevantes y la estrategia de pruebas.
\item
  \textbf{Capítulo 5 --- Resultados y Validación}: presenta los resultados de las pruebas de rendimiento, usabilidad y validación con empresas piloto, contrastándolos con los objetivos planteados.
\item
  \textbf{Capítulo 6 --- Conclusiones y Trabajo Futuro}: sintetiza los logros del proyecto por objetivo, enumera las limitaciones encontradas y propone líneas de trabajo para versiones futuras.
\item
  \textbf{Referencias Bibliográficas}: lista completa de las fuentes citadas a lo largo del documento.
\item
  \textbf{Anexos}: manual de usuario, manual de despliegue, especificación de la API y resultados detallados de pruebas.
\end{itemize}
